% ============================================================
%  Chapitre 1 — Contexte général
% ============================================================
\chapter{Contexte général}
\label{ch:contexte}

\section{Introduction}

Ce chapitre situe le cadre dans lequel le projet a été réalisé. Nous présentons
d'abord l'organisme d'accueil, sa vision et ses valeurs. Nous exposons ensuite la
méthodologie de conduite de projet retenue --- CRISP-DM --- en détaillant ses
six phases et en précisant la manière dont chacune s'applique à un projet de
type RAG, qui présente la particularité de ne pas reposer sur l'entraînement
d'un modèle mais sur l'assemblage et le réglage d'une chaîne de composants.

\section{Présentation de l'organisme d'accueil}

Créée en 2010, \organism{} est la structure dédiée à la Recherche, au
Développement et à l'Innovation (RDI) au sein de l'école ESPRIT. Depuis la mise
en place des programmes de formation, l'activité RDI est devenue l'une des
priorités stratégiques de l'établissement.

\organism{} regroupe des enseignants-chercheurs et offre un soutien complet en
matière de financement, de formation et d'accompagnement dans les projets de
recherche. La structure définit les stratégies de recherche en collaboration
avec les différentes parties prenantes, en mettant particulièrement l'accent sur
l'innovation et la recherche appliquée. C'est dans ce cadre que le présent
projet a été mené, à l'intersection du traitement automatique du langage naturel
et de l'ingénierie logicielle.

\section{Cadre stratégique et valeurs de l'organisme}

\organism{} a pour vision de promouvoir la recherche, l'innovation et la
valorisation des résultats scientifiques. Elle s'engage à soutenir l'excellence
en RDI à travers le développement de projets à fort impact, répondant aux
besoins académiques et industriels.

Sa mission consiste à créer un environnement propice à la recherche et à
l'innovation, à renforcer le financement externe et à développer un écosystème
entrepreneurial ouvert à l'international, favorisant la collaboration entre
chercheurs et partenaires.

Les valeurs de \organism{} reposent sur la rigueur scientifique, le respect,
l'inclusion et la transparence. Cette exigence de transparence a d'ailleurs eu
une traduction technique directe dans le présent projet : le système développé
ne se contente pas de produire une réponse, il expose systématiquement les
extraits documentaires sur lesquels cette réponse s'appuie, avec leur score de
pertinence. L'utilisateur peut ainsi vérifier ce que le système avance plutôt
que de devoir lui faire confiance.

\section{Méthodologie CRISP-DM}

CRISP-DM, ou \eng{Cross-Industry Standard Process for Data Mining}, est un modèle
de processus qui permet d'orienter les travaux d'exploration de données en
science des données. Cette méthodologie itérative permet de travailler en
cycles : de la compréhension du besoin métier jusqu'au déploiement, en passant
par la préparation des données, la modélisation et l'évaluation. Elle donne un
cadre structuré pour guider le \eng{data scientist} tout au long du cycle de vie
d'un projet.

Son choix pour ce stage se justifie par deux caractéristiques. D'une part, son
caractère \textbf{itératif} : dans un projet RAG, la qualité finale dépend de
paramètres --- taille des segments, nombre de documents récupérés, formulation
de l'invite --- dont la valeur optimale ne peut être déterminée
\emph{a priori}. Il faut mesurer, ajuster, remesurer. D'autre part, la place
centrale qu'elle accorde à la phase d'\textbf{évaluation} avant tout
déploiement : c'est précisément le point critique d'un système génératif, dont
la sortie est du texte en langage naturel et non une classe ou une valeur
numérique aisément comparable à une vérité de terrain.

\section{Les étapes de la méthodologie CRISP-DM}

La méthodologie CRISP-DM se compose de six étapes clés.

\subsection*{Étape 1 : la compréhension du besoin métier}

Dans cette phase initiale, on travaille en collaboration avec les parties
prenantes afin de comprendre les objectifs, les exigences et les contraintes du
projet ainsi que les bénéfices attendus. Les ressources nécessaires --- budget,
compétences techniques, accès aux données --- sont évaluées. Les problèmes à
résoudre sont identifiés et les critères permettant de mesurer le succès du
projet sont définis. Pour ce projet, cette étape a produit la contrainte
structurante du « tout local » et le critère de succès central : la traçabilité
de chaque réponse. Elle fait l'objet du chapitre~\ref{ch:metier}.

\subsection*{Étape 2 : la compréhension des données}

Il s'agit de collecter les données disponibles, d'en décrire la structure, d'en
évaluer la qualité et d'identifier les premières tendances. Dans un projet RAG,
cette étape porte sur le corpus documentaire : quels types de documents, quel
volume, quelle qualité d'extraction textuelle, quelle structure interne. Elle
détermine directement la stratégie de découpage adoptée ensuite.

\subsection*{Étape 3 : la préparation des données}

Phase généralement la plus coûteuse en temps, elle consiste à nettoyer,
transformer et mettre en forme les données pour l'étape suivante. Ici, elle
recouvre l'extraction du texte des fichiers PDF, l'enrichissement en métadonnées
de traçabilité et le découpage en segments de taille compatible avec le modèle
de vectorisation. Les étapes 2 et 3 constituent le chapitre~\ref{ch:donnees}.

\subsection*{Étape 4 : la modélisation}

On construit des modèles prédictifs ou descriptifs en fonction des objectifs du
projet. Plusieurs itérations ont lieu entre les étapes de préparation et de
modélisation pour affiner l'usage de certains algorithmes.

Cette étape mérite une remarque particulière dans le cadre d'un projet RAG.
Contrairement à un projet d'apprentissage supervisé classique, aucun modèle
n'est ici \emph{entraîné} : le modèle de vectorisation et le modèle de langue
sont tous deux pré-entraînés et utilisés en inférence. La « modélisation »
consiste à \textbf{assembler et paramétrer une chaîne de traitement} :
choix du modèle d'\eng{embeddings}, structure de l'index vectoriel, stratégie de
recherche, formulation de l'invite système. Ce sont ces choix, et non des poids
appris, qui déterminent la performance du système. C'est l'objet du
chapitre~\ref{ch:modelisation}.

\subsection*{Étape 5 : l'évaluation}

Une fois les modèles construits, ils sont évalués afin de déterminer leur
qualité. Cette étape permet de s'assurer que le système atteint les objectifs du
projet. Les performances sont mesurées à l'aide de métriques appropriées et
comparées aux critères de succès définis en première phase. Si les résultats ne
répondent pas aux attentes, les étapes précédentes peuvent être révisées et
répétées.

Pour un système génératif, l'évaluation pose une difficulté propre : il n'existe
pas de métrique unique comme l'\eng{accuracy}. Nous verrons au
chapitre~\ref{ch:evaluation} comment le \eng{framework} RAGAS décompose la
qualité d'un système RAG en quatre dimensions mesurables séparément, dont une
--- la \eng{faithfulness} --- mesure directement le taux d'hallucination.

\subsection*{Étape 6 : le déploiement}

Enfin, les résultats du projet sont présentés aux parties prenantes et intégrés
si nécessaire dans les systèmes existants pour aider à la prise de décision.
Cette phase implique la création de rapports, de visualisations ou d'autres
formes de communication rendant les résultats compréhensibles et utilisables par
des non-spécialistes. Elle correspond au chapitre~\ref{ch:deploiement} : API
REST, interface conversationnelle et conteneurisation.

\vspace{0.5cm}

La figure~\ref{fig:crispdm} illustre le caractère cyclique de ce processus.

\begin{figure}[htbp]
\centering
\begin{tikzpicture}[
  node distance=1.05cm and 1.5cm,
  phase/.style={
    draw=accent,line width=1pt,rounded corners=5pt,fill=soft,
    text width=3.1cm,align=center,minimum height=1.05cm,font=\small\bfseries},
  data/.style={
    draw=ink,line width=1pt,circle,fill=white,
    text width=1.7cm,align=center,font=\small\bfseries},
  fl/.style={-{Stealth[length=2.6mm]},line width=0.9pt,color=muted}
]

\node[data] (data) {Données};

\node[phase,above left=0.9cm and 0.55cm of data]  (bu) {Compréhension\\métier};
\node[phase,above right=0.9cm and 0.55cm of data] (du) {Compréhension\\des données};
\node[phase,right=2.05cm of data]                 (dp) {Préparation\\des données};
\node[phase,below right=0.9cm and 0.55cm of data] (mo) {Modélisation};
\node[phase,below left=0.9cm and 0.55cm of data]  (ev) {Évaluation};
\node[phase,left=2.05cm of data]                  (de) {Déploiement};

\draw[fl] (bu) -- (du);
\draw[fl] (du) -- (dp);
\draw[fl] (dp) -- (mo);
\draw[fl] (mo) -- (ev);
\draw[fl] (ev) -- (de);
\draw[fl] (de) -- (bu);

% aller-retours internes caractéristiques de CRISP-DM
\draw[fl,dashed] (du) to[bend right=18] (bu);
\draw[fl,dashed] (mo) to[bend left=18]  (dp);
\draw[fl,dashed] (ev) to[bend left=22]  (bu);

\end{tikzpicture}
\caption{Le processus CRISP-DM et ses six phases itératives. Les flèches en
pointillés représentent les retours en arrière fréquents, notamment entre
modélisation et préparation des données.}
\label{fig:crispdm}
\end{figure}

\section{Conclusion}

Ce chapitre a présenté \organism{}, son rôle, sa vision et ses valeurs, ainsi
que la méthodologie CRISP-DM retenue pour structurer le projet. Nous avons
souligné les deux adaptations que requiert cette méthodologie dans un contexte
RAG : la modélisation devient un travail d'assemblage et de paramétrage plutôt
que d'entraînement, et l'évaluation exige des métriques spécifiques aux systèmes
génératifs. Le chapitre suivant formalise la problématique métier à laquelle le
système doit répondre.
